% MCV-centered revised introduction for Overleaf.
% Citations are written as author-year placeholders for later integration.

\documentclass[11pt]{article}
\usepackage[margin=1in]{geometry}
\usepackage{amsmath}
\usepackage[version=4]{mhchem}
\usepackage[T1]{fontenc}
\usepackage{lmodern}

\title{Precession-Phase Organization of DO-Like Asian Monsoon Transitions}
\author{}
\date{}

\begin{document}
\maketitle

\section*{Introduction}

Millennial-scale climate variability (MCV) has punctuated glacial climate throughout much of the Pleistocene. Its best-known expression is the Dansgaard--Oeschger (DO) sequence of the last glacial period, in which Greenland temperatures shifted abruptly between stadial and interstadial states on decadal to centennial timescales (Rasmussen et al., 2014; Kindler et al., 2014; Capron et al., 2021). These abrupt reorganizations were not confined to Greenland. They were associated with widespread changes in atmospheric circulation, monsoon hydroclimate, the global water cycle, and interhemispheric heat redistribution (Clement et al., 2001; Chiang and Bitz, 2005; Corrick et al., 2020). The recurrence, abruptness, and broad climatic imprint of MCV make it a central problem for understanding the stability of glacial climate states.

The immediate expression of many DO-style transitions is commonly linked to reorganizations of the Atlantic Meridional Overturning Circulation (AMOC), but the boundary conditions that control the occurrence, strength, and pacing of MCV remain debated (Broecker, 1998; Clark et al., 2002; Zhang et al., 2014). A growing body of evidence indicates that orbital forcing and glacial background state modulate MCV amplitude. Sun et al. (2021) showed that MCV persisted over the past \(\sim1.5\) Myr and that its amplitude was influenced by precession, obliquity, and glacial boundary conditions. North Atlantic records likewise suggest coupling between orbital and millennial variability, with MCV amplitude varying with ice volume and orbital configuration (Hodell et al., 2023). In the monsoon domain, Thirumalai et al. (2020) used the Cheng et al. (2016) Chinese speleothem composite, together with atmospheric \(\ce{CH4}\), to examine orbital modulation of millennial-scale variability. Their results placed Chinese speleothem millennial variability within the broader MCV proxy framework and showed that orbital forcing can modulate the magnitude of millennial-scale signals, while also highlighting proxy-dependent expressions of that modulation.

Amplitude modulation, however, does not necessarily imply event pacing. A larger millennial-scale envelope under a particular orbital configuration could arise because externally forced boundary conditions amplify internally generated transitions, without determining when individual transitions occur. In stochastic or excitable climate systems, abrupt regime shifts can be triggered when internal variability crosses a threshold, producing clustered or intermittent events even without direct external pacing (Ditlevsen and Ditlevsen, 2005; Kleppin et al., 2015; Lohmann and Ditlevsen, 2018). Conversely, climate-model experiments suggest that orbital configuration can directly affect DO-like oscillatory regimes and abrupt transitions. For example, changes in precession or obliquity can alter tropical hydroclimate, high-latitude sea ice, and ocean circulation in ways that switch millennial-scale AMOC oscillations on or off under otherwise fixed boundary conditions (Zhang et al., 2021). The key unresolved proxy question is therefore not only whether orbital forcing modulates MCV amplitude, but whether it also organizes the timing of individual stadial--interstadial-like transitions.

Testing this event-timing question requires a long and precisely dated proxy sequence of abrupt transitions. Greenland ice-core chronologies provide the canonical DO sequence, but their temporal coverage is largely limited to the last glacial cycle, which restricts the separation of precession, obliquity, and glacial-background effects. Chinese speleothem records provide a complementary archive. During the last glacial period, independently dated speleothem records show that abrupt Greenland warmings were accompanied by near-synchronous hydroclimate shifts in monsoon regions (Corrick et al., 2020), supporting the use of speleothem abrupt monsoon changes as regional proxies for MCV. The 640 kyr Chinese speleothem composite of Cheng et al. (2016), and the weak and strong monsoon starts objectively identified from it by Rousseau et al. (2023), therefore offer a rare opportunity to test event timing over multiple orbital cycles. These weak and strong starts are not Greenland-defined DO events themselves; rather, they are DO-like Asian monsoon transitions that record the monsoon-sector expression of millennial-scale climate variability.

Here we test whether the occurrence timing of these DO-like Asian monsoon transitions is organized by orbital phase. We use the 0--640 ka Rousseau et al. (2023) weak and strong monsoon-start catalogues and ask three linked questions. First, are the event sequences consistent with stationary random occurrence? Second, do weak and strong monsoon starts cluster at preferred precession or obliquity phases? Third, does any orbital-phase dependence remain after accounting for slow background states represented by LR04 ice volume and atmospheric \(\ce{CO2}\)? To answer these questions, we combine Lohmann-style stationary randomness tests, circular Rayleigh phase tests, and binned Poisson hazard generalized linear models. This event-based framework extends previous work on orbital modulation of MCV amplitude by asking whether orbital configuration also affects transition probability.

We find that weak and strong monsoon starts are inconsistent with simple stationary randomness and exhibit significant, opposite preferences in precession phase. Poisson hazard models further show that precession phase remains informative after accounting for LR04 and \(\ce{CO2}\), whereas obliquity phase and other candidate forcings do not provide comparable independent contributions. These results suggest that precession acted not only as a modulator of millennial-scale variability amplitude, but also as a phase-dependent boundary condition for the occurrence probability of DO-like Asian monsoon transitions.

\end{document}
